\chapter{Ergebnisse aus der Messung}
\section{Datengrundlage}
Die Datengrundlage für die Auswertung des stationären EEG-Systems BioSemi
ActiveTwo bilden die Daten aller zehn Versuchspersonen. Für das portable
EEG-System von OpenBCI liegen nur neun auswertbare Datensätze aus
Versuchsbedingung 1 (außerhalb der Kabine) und sechs aus Versuchsbedingung 2
(innerhalb der Kabine) vor. Grund dafür ist das Ausschließen einiger
Datensätze, aufgrund von Elektroden, dessen Signale keine ausreichende Qualität
aufwiesen. Eine Ursache sind die Trockenelektroden des OpenBCI Systems, welche
bei einigen Messungen keinen dauerhaft ausreichenden Kontakt zur Kopfhaut
hatten. Die Signale der betroffenen Elektroden waren dadurch entweder stark
verrauscht oder übersteuert, das heißt, sie zeigten dauerhaft den maximal
messbaren Wert an und enthielten keine verwertbare Hirnaktivität mehr. Solch
übersteuerte Elektroden werden auch als \textit{railed} bezeichnet. Einzelne
gestörte Elektroden lassen sich zwar aus den benachbarten Elektroden
rechnerisch durch Interpolation ersetzen (vgl. Abschnitt~3.4.1), sind jedoch
zu viele Elektroden betroffen, ist dieses Verfahren nicht anwendbar, da die
Interpolation fehlender Kanäle aus benachbarten Elektroden voraussetzt, dass
genügend gute Nachbarkanäle vorhanden sind. Folgende Datensätze mussten
dadurch ausgeschlossen werden: in Versuchsbedingung 1 der Datensatz von
\gls{pt}\_06 und in Bedingung 2 die Datensätze von pt\_01, pt\_04,
pt\_07 und pt\_06. \newline 
Die Hirnaktivität unterscheidet sich zwischen einzelnen Personen deutlich,
beispielsweise in der Stärke der Alpha-Aktivität oder in der Amplitude der P1. Würde die Gruppenmittelungen der einzelnen Messungen mit unterschiedlichen
Versuchspersonen gebilden werden, könnte ein Unterschied zwischen den Messungen allein
darauf zurückzuführen sein, dass jeweils andere Personen in die Mittelung eingehen, und
nicht auf dem verwendeten EEG-System oder der Versuchsbedingung. Deshalb wurden für den Vergleich der Messungen nur Versuchspersonen herangezogen, die 
die in beiden Messungen gültige Daten aufweisen. 

\section{Verhaltensergebnisse}
Zur Erfassung der Verhaltensdaten wurden pro Versuchsperson im Kartendeck acht
Könige gezeigt, auf welche mit dem Drücken der Leertaste reagiert werden
sollte. Dies sollte die Aufmerksamkeit der Proband/innen sicherstellen. Als
Reaktionszeit wurde die Zeit zwischen dem Erscheinen eines Königs und dem
Drücken der Leertaste bestimmt. Anschließend wurde der Median der
Reaktionszeiten berechnet und dieser Wert über alle Versuchspersonen gemittelt.

Die Aufgabe wurde in allen Messungen nahezu fehlerfrei gelöst, sprich auf alle
Könige wurde mit Drücken der Leertaste reagiert. Von 80 Königen pro Messung
wurden in Bedingung~1 und bei BioSemi alle, in Bedingung~2 79 Könige erkannt.
Die Reaktionszeit (Mittelwert $\pm$ Standardabweichung der individuellen
Mediane) betrug in Bedingung~1 $550 \pm 68$\,ms, in Bedingung~2
$589 \pm 79$\,ms und bei BioSemi $591 \pm 70$\,ms. Die Messungen
unterschieden sich damit in der Reaktionszeit. In Bedingung~2 reagierten alle
zehn Versuchspersonen langsamer als in Bedingung~1, während sich Bedingung~2
und BioSemi bezüglich des Verhaltens nicht unterschieden. Ordnet man die
Messungen nach ihrer Position im zeitlichen Abfolge, war bei allen zehn
Versuchspersonen die jeweils erste Messung die schnellste. Die Reaktionszeit
verlängert sich also nach der ersten Messung um etwa 40\,ms und bleibt danach
konstant.
\section{Ergebnisse der ERP-Analyse}
Ausgewertet wurde die mittlere P1-Amplitude im Zeitfenster von 80--200\,ms an den okzipitalen Elektroden O1 und O2. Das Fenster wurde gewählt, da die Latenz der P1 je nach Paradigma und Stimuluseigenschaften variieren kann und der P1-Gipfel in den vorliegenden Daten später auftrat als in der Literatur beschrieben (vgl. Abschnitt~\ref{sec:erp}). Um zu prüfen, ob sich die P1 zwischen den Bedingungen unterscheidet, wurde eine \gls{varianzanalyse} mit \gls{messwiederholung} berechnet.

Für das BioSemi-System wurde untersucht, ob die P1 von der Darbietungsform (groß und normalkontrastiert oder klein und kontrastreduziert) abhängt. Für das OpenBCI-System wurde zusätzlich geprüft, ob die Versuchsbedingung (Kabine geöffnet oder geschlossen) einen Einfluss auf die P1 hat. Anschließend wurden einzelne Bedingungen mit gepaarten t-Tests direkt miteinander verglichen, zum Beispiel große mit kleinen Stimuli. Im Folgenden werden nur die Ergebnisse berichtet, die für die Fragestellung relevant sind. Nicht berichtet wird, ob die P1 insgesamt auf der linken oder rechten Kopfseite größer war.

\begin{table}[htbp]
  \centering
  \caption{Mittlere P1-Amplitude (80--200\,ms) in \textmu V, Mittelwert $\pm$ Standardabweichung}
  \label{tab:p1_amplituden}
  \begin{tabular}{llccc}
    \hline
    System & Versuchsbedingung & $n$ & groß, normalkontrastiert & klein, kontrastreduziert \\
    \hline
    BioSemi & VB 3 & 10 & 4,70 $\pm$ 3,19 & 3,91 $\pm$ 1,91 \\
    OpenBCI & VB 2 (Kabine geschlossen) & 6 & 2,37 $\pm$ 0,85 & 1,99 $\pm$ 1,10 \\
    OpenBCI & VB 1 (Kabine geöffnet) & 6 & 1,35 $\pm$ 1,11 & 1,20 $\pm$ 1,41 \\
    \hline
  \end{tabular}
\end{table}

Tabelle~\ref{tab:p1_amplituden} zeigt die mittleren P1-Amplituden für beide Darbietungsformen in allen drei Versuchsbedingungen. 


\section{Ergebnisse der Frequenzbandanalyse}
\section{Gegenüberstellung der Daten anhand von Performanzkriterien}
